\documentclass[accentcolor=tud1b,colorbacktitle,inverttitle,landscape,presentation,t]{tudbeamer}
\usepackage[utf8]{inputenc}
\usepackage[german]{babel}
\usepackage{verbatim, hyperref,multimedia}
\usepackage{graphicx, psfrag}			%to insert graphics
\usepackage{setspace}
\usepackage[active]{srcltx}
\usepackage{color}
\usepackage{subfigure}
\usepackage{listings}
\usepackage{pstricks}
\usepackage{amssymb}

\useinnertheme{default}
\usecolortheme{orchid}
\setbeamercovered{transparent}


\newcommand{\keyword}[1]{\textcolor{tudaccent}{\textbf{#1}}} 
\newcommand{\myframetitle}[2]{\frametitle{#1 \\[.2cm] \small #2}}
\newcommand{\fatitem}[1]{\item \textbf{#1}}
\newcommand{\tab}{\hspace*{0.5cm}}


\setcounter{tocdepth}{1} 

\hypersetup{
pdftitle={WebMining Exercise }, pdfauthor={F. Englert}
}


\begin{document}

\title[MGA]{\large 4. Exercise}

\author{Frank Englert, Jens Haase}
\institute{KE, TU Darmstadt}

\logo{\includegraphics{graphics/keico}}

\date{May 2011}

\begin{titleframe}
\begin{center}
\color{tudtextaccent} \large Die 4. Übung\\[.5cm]
%\includegraphics[width=0.4\textwidth]{Graphics/spg_logo_final.eps} \\
\normalcolor \normalsize Knowlege Engineering \\
Fachbereich Informatik \\
Technische Universität Darmstadt\\[.5cm]

\textbf{Präsentation der Ergebnisse:}\\
Frank Englert\\
Jens Haase
\end{center}

\end{titleframe} 

\begin{frame}[c]
	\myframetitle{4. Übung}{Inhalt}
\begin{enumerate}
  \item Datenaufbereitung
  \item ​Sim­ple Naive Bayes
  \item Compress Al­go­rith­mus
  \item Eval­u­a­tion
\end{itemize}
\end{enumerate}
\end{frame}

\begin{frame}[c]
	\myframetitle{Task 1}{Datenerfassung und Aufbereitung}
\begin{itemize}
  \fatitem{Verwendung des WebKB-Datensatzes}
  \item Aufteilung in Traininings- und Testdaten im Verhältnis von 1/1
  \item Daraus resultieren ~4k Testdaten
  \item Klassifikation nach folgenden Klassen:
  \begin{itemize}
  \item course
  \item department
  \item faculty
  \item other
  \item project
  \item stuff
  \item student
\end{itemize}
 
\end{itemize}
\end{frame}


\begin{frame}[c]
	\myframetitle{Task 2.1 }{Textextraktion}
\begin{itemize}
  \fatitem{ Textextraktion erfolgt durch den JSoup HTML-Parser}
  \begin{itemize}
  \item Html-Tags werden bis auf den Text-Content entfernt
  \item Beispiel: <p> Webmining Übung <i>3</i> $\rightarrow$ Webmining Übung 3 
\end{itemize}
  \fatitem{Wortextraktion erfolgt durch die Klasse BreakIteratorTokenizer}
  \begin{itemize}
  \item Diese ist Teil des Apache UIMA Projekts
  \fatitem{Zum Stemmen wird der Snowball-Stemmer verwendet}
  \item \end{itemize}
\end{itemize}
\end{frame}
\begin{frame}[c]
	\myframetitle{Task 3}{Compress}
\begin{itemize}
  \item ...
 
\end{itemize}
\end{frame}


\begin{frame}[c]
	\myframetitle{Task 4}{Evaluation}
\begin{itemize}
  \item ...
 
\end{itemize}
\end{frame}



\end{document}
